\begin{abstract}
Post-hoc model merging has emerged as a groundbreaking, training-free paradigm for multi-task learning, enabling separate expert networks independently fine-tuned from a shared initialization to be directly composed in weight space. To mitigate task conflicts, recent frameworks employ test-time adaptation (TTA) to dynamically optimize layer-wise merging coefficients. However, we identify a critical, overlooked vulnerability in these adaptive paradigms: \textbf{Quantization-Operator Overfitting}. Unregularized coefficient optimization converges to extremely sharp local minima that are highly fragile under weight perturbations, triggering catastrophic performance collapse when the merged model is subjected to downstream deployment post-training quantization (PTQ). 

To resolve this bottleneck, we propose \textbf{Hessian-Regularized Coefficient Optimization (HessMerge)}, a mathematically rigorous framework that guarantees PTQ robustness. By analyzing the quantization process as a parameter-space perturbation, we show that the quantization-induced loss gap is bounded by the curvature of the loss landscape. HessMerge explicitly incorporates a regularizer penalizing the \textbf{trace of the Hessian} ($\text{Tr}(\mathcal{H}_{\Theta})$) of the multi-task loss with respect to the continuous merging coefficients, flattening the local loss minimum during test-time adaptation. We evaluate HessMerge against state-of-the-art baselines under six diverse quantization schemas (ranging from FP32 down to INT4 symmetric per-channel) across four visual domain datasets (MNIST, FashionMNIST, CIFAR-10, SVHN) using a Vision Transformer backbone. Our empirical findings demonstrate that HessMerge dramatically improves robustness to quantization noise, maintaining high joint multi-task accuracy even under aggressive INT4 quantization and establishing a new state of the art in robust model merging.
\end{abstract}
